\section{The logit taste scale: smoothing device and structural elasticity}
\label{app:taste}

The extreme-value taste shocks on the adoption choice play two roles at once, and
this appendix separates them. The first role is computational: the logit smooths
the discrete policy, so that choice probabilities and the value function are
differentiable with closed-form derivatives, which is what the sequence-space
method requires to build Jacobians through the household block
\citep{DCEGM2017}. In that tradition the scale $\sigma_\varepsilon$ is a
numerical device: one takes it small and verifies that results are insensitive to
it. The second role is economic. Differentiating the logit gives
$dP_i = P_i\,(dV_i - dV)/\sigma_\varepsilon$: the response of the choice
probabilities to \emph{any} perturbation of the value gap scales with
$1/\sigma_\varepsilon$, so the same parameter that smooths the policy is the
structural elasticity of the margin. Which of the two readings governs is not a
matter of taste; it depends on where the population sits on the logit.

If choice probabilities are interior, of order one, as for labor-force
participation, where employment rates and quarterly flows are tens of
percent, then $P(1-P)$ is bounded, the $\sigma_\varepsilon \to 0$ limit is a
well-defined deterministic cutoff, and aggregate responses converge to the
density of households at the threshold: an object of the wealth and productivity
distribution, not of $\sigma_\varepsilon$. There the numerical reading is
legitimate, the economics of the margin is carried by the distribution (and, in
search applications, by the frictions standing between the choice and the
outcome), and the smoothing doctrine applies. If instead the margin is
\emph{rare}, choice probabilities deep in the left tail, the logit is
exponential there, $d\log P = d\Delta/\sigma_\varepsilon$ holds for every
household regardless of its level, and the semi-elasticity of the flow is
identically $1/\sigma_\varepsilon$. No feature of the distribution can attenuate
it, there is no $\sigma_\varepsilon$-insensitive range, and the scale \emph{is}
the elasticity of the margin.

Our margin is rare, verifiably so. The calibration targets a $5\%$ green share
with a five-year durable life, which fixes the steady-state switching flow at
$\delta_g \bar D^{G}/(1-\bar D^{G}) = 0.26\%$ per quarter.
Table~\ref{tab:switch_prob_dist} reports the distribution-weighted quantiles of
the quarterly switch probability among brown incumbents: the mean is $0.26\%$,
the median is essentially zero, the 95th percentile is $1.6\%$, and even the
99.9th percentile stays below $9\%$. Virtually no mass approaches interior
probabilities. This is the world of rare, lumpy, quasi-irreversible
adjustment, the regime in which the structural discrete-choice literature
estimates the scale from choice data \citep{Rust1987},not the world of
participation.

\input{output/tab_switch_prob_dist_import.tex}

Three consequences follow, each handled in the paper. First, the adoption-channel
magnitudes are proportional to $1/\sigma_\varepsilon$, so a level moment cannot
discipline them; Section~\ref{sec:taste_id} characterises the
$(\sigma_\varepsilon,\psi_g)$ locus and identifies the adoption elasticity as the
moment that pins the scale. Second, the same $1/\sigma_\varepsilon$ governs the
curvature of the margin, so the linearisation error is concentrated in $D^{G}$
and grows as $\sigma_\varepsilon$ falls; Appendix~\ref{sec:nonlin} quantifies it.
Third, the steady state and the policy ordering are invariant along the locus, so
none of the paper's qualitative results turn on the scale. What does turn on
it, the size of the adoption wave and of the adoption channel, is exactly what
the elasticity moment disciplines.

\subsection{From micro estimates to the taste scale: the procedure}
\label{app:taste_procedure}

This subsection documents the calibration of $\sigma_\varepsilon$ step by step,
so that every number in Section~\ref{sec:taste_id} can be audited and
reproduced. The entire procedure is implemented in a single runner
(\texttt{runners/run\_taste\_identification.py}; \texttt{python main.py
taste\_id} regenerates Table~\ref{tab:taste_identification},
Table~\ref{tab:switch_prob_dist} and Figure~\ref{fig:taste_identification});
the grid, the perturbation sizes and the dollar mapping are constants at the top
of the file.

\begin{enumerate}\itemsep3pt
\item \emph{Trace the locus.} For each $\sigma_\varepsilon$ on the grid
$\{0.02, 0.035, 0.05, 0.07, 0.10, 0.15, 0.20\}$, re-solve the baseline steady
state with $\psi_g$ as the unknown targeting $\bar D^G = 5\%$. This delivers
$\psi_g^*(\sigma_\varepsilon)$, rising from $0.32$ to $1.53$, with the steady
state, shares, flows, prices, distribution, invariant along the locus.

\item \emph{Switch to the fixed-$\psi$ configuration.} All comparative statics
below hold $\psi_g$ fixed at $\psi_g^*$ and let the green share $D^G$ adjust
endogenously (the same configuration that builds the ex-ante ETS economy of
Section~\ref{sec:versionA}). Elasticities are computed as finite differences on
this configuration around the calibrated point.

\item \emph{Operating-cost margin.} Perturb the delivered green/brown price
ratio $\varrho = 0.80$ down by $0.01$ at fixed $\psi_g^*$ and measure
$\Delta \ln D^G / \Delta \ln (1-\varrho)$, the elasticity of the green share to
the operating-cost gap. It falls from $2.8$ to $0.4$ across the grid ($1.32$ at
the baseline).

\item \emph{Upfront-cost margin: the matched moment.} Perturb $\psi_g$ down by
$1\%$ at fixed calibration and measure the response of the switching flow
$P = D^{sw}/(1-D^G)$. Per $1\%$ of $\psi_g$ the semi-elasticity is nearly flat
in $\sigma_\varepsilon$ ($5.1\%$ to $3.8\%$), because $\psi_g^*$ rises along the
locus and compensates; the identifying variation comes from converting to
absolute dollars. With one model unit equal to quarterly household consumption
$\bar C_{\$}$, the semi-elasticity per \$1{,}000 is the per-$1\%$ figure times
$1{,}000/(0.01\,\psi_g^*\,\bar C_{\$})$; at $\bar C_{\$} = \$20{,}000$ it falls
from $80\%$ to $12\%$ across the grid
(Table~\ref{tab:taste_identification}).
\end{enumerate}

The empirical anchors are converted to the same units.
\citet{GallagherMuehlegger2011} and \citet{Chandra2010} estimate that a
\$1{,}000 tax incentive raises hybrid adoption by $31$--$38\%$; no conversion is
needed. \citet{MuehleggerRapson2022} report a demand elasticity of $-2.1$ at a
mean transaction price of \$26{,}000, which converts to
$2.1 \times 1{,}000/26{,}000 \approx 8\%$ per \$1{,}000; their estimated
pass-through of $73$--$85\%$ implies the price faced by the buyer moves nearly
one-for-one with the subsidy, so the conversion needs no incidence correction.
The choice between the two anchors is a population statement, not a taste: at a
$5\%$ green share the model's marginal adopters consume roughly twice the
average brown household, they are the higher-income early adopters of the
hybrid-era estimates, not the mass-market population of
\citet{MuehleggerRapson2022}, whose estimate therefore serves as the floor on
responsiveness rather than the target.

Matching the early-adopter range pins $\sigma_\varepsilon \in [0.05, 0.07]$. The
implied premium $\psi_g^* \bar C_{\$}$ provides an independent cross-check: it is
\$11{,}000--\$14{,}000 on the pinned interval, within the heat-pump/EV premium
range, while $\sigma_\varepsilon \geq 0.10$ would imply premia of \$18{,}000 and
above.

The one free constant is the dollar mapping $\bar C_{\$}$, an assumption stated,
not estimated. Its role is transparent: a higher $\bar C_{\$}$ scales down the
model's per-\$1{,}000 response and shifts the elasticity-based pin toward lower
$\sigma_\varepsilon$, while scaling up the implied premium and shifting the
premium-based cap in the same direction, so the two criteria co-move and the
joint interval stays tight. Concretely, at $\bar C_{\$} = \$16{,}000$ the joint
pin is roughly $[0.07, 0.10]$; at \$20{,}000, $[0.05, 0.07]$; at \$24{,}000,
$[0.04, 0.06]$. Every one of these intervals lies inside the sensitivity grid of
Table~\ref{tab:taste_identification}, over which the steady state and the policy
ordering are invariant; what the mapping moves is the size of the adoption
channel, bounded on the grid by peak responses of $2$ to $20$ pp.

Two limits of the procedure are worth stating. The anchors come from vehicle
adoption in the United States in the 2000s--2010s; their external validity for
heat pumps or for European households is an assumption, and estimating the
adoption elasticity directly, rather than importing it, is the natural next
step (Section~\ref{sec:limits}). And the anchors identify the upfront-cost
margin; the operating-cost margin, the channel documented by
\citet{BeresteanuLi2011} for gasoline prices and hybrid demand, is disciplined
here only through the model's internal consistency, which ties the two margins
together at the calibrated point.

\begin{figure}[H]
\centering
\includegraphics[width=\textwidth]{output/fig_taste_identification.png}
\caption{Identification of the taste scale. \emph{Left}: the steady-state adoption
elasticity (green share vs.\ the operating-cost gap) against $\sigma_\varepsilon$.
\emph{Right}: the crisis adoption magnitude (peak $\Delta D^G$ to the brown-price
shock). The steady state is held at $\bar D^G=5\%$ by recalibrating $\psi_g$; the
dashed line marks the baseline $\sigma_\varepsilon=0.05$.}
\label{fig:taste_identification}
\end{figure}

\begin{table}[H]
\centering\small
\begin{tabular}{lrrrrrr}
\toprule
$\sigma_\varepsilon$ & $\psi_g$ & premium \$000 & cost-gap elast. & \%/\$1{,}000 & peak $\Delta D^G$ & $\sum y$ \\
\midrule
$0.020$ & $0.323$ & $6.5$ & $2.78$ & $80$ & $19.86$ & $-59.2$ \\
$0.035$ & $0.430$ & $8.6$ & $1.80$ & $54$ & $12.42$ & $-60.9$ \\
$0.050$ & $0.538$ & $10.8$ & $1.32$ & $40$ & $8.87$ & $-61.5$ \\
$0.070$ & $0.680$ & $13.6$ & $0.97$ & $31$ & $6.34$ & $-61.6$ \\
$0.100$ & $0.888$ & $17.8$ & $0.69$ & $22$ & $4.36$ & $-61.4$ \\
$0.150$ & $1.220$ & $24.4$ & $0.47$ & $16$ & $2.80$ & $-60.6$ \\
$0.200$ & $1.534$ & $30.7$ & $0.36$ & $12$ & $2.01$ & $-59.8$ \\
\bottomrule
\end{tabular}
\caption{Identification of the logit taste scale. For each $\sigma_\varepsilon$,
$\psi_g$ is recalibrated to hold $\bar D^G=5\%$, leaving the steady state invariant.
Columns: the implied green premium in dollars (one unit $=\$20{,}000$ quarterly
consumption); the steady-state elasticity of the green share to the operating-cost
gap at fixed $\psi_g$; the semi-elasticity of the switching flow to a \$1{,}000 cut
in $\psi_g$, matched to \citet{GallagherMuehlegger2011,Chandra2010,MuehleggerRapson2022}; the peak
crisis adoption response; and the cumulative output loss, flat while the adoption
responsiveness varies by an order of magnitude.}
\label{tab:taste_identification}
\end{table}
